\documentclass[11pt,english,francais]{article}
\usepackage{scrtime}
\usepackage{natbib} 
\usepackage{tikz} 
\usepackage{multirow}
\usepackage[utf8]{inputenc}
\usepackage{amsmath} 
\usepackage{amsfonts} 
\usepackage{amsthm} 
\usepackage{thmtools}
\usepackage{hyperref}
\usepackage{cleveref}
\usepackage[hypcap]{caption}
\usepackage[off]{auto-pst-pdf}
\usepackage[scale=2]{ccicons}
\usepackage{tabularx}
\DeclareMathOperator{\sgn}{sgn}

\begin{document}


\section{Méthodes de lissage et saisonnalité}
\label{sec:serie-dexercices-1}

\subsection{Noël s'en vient !}
\label{sec:exercice-1-1}

On considère les taux d'inflation sur 12 mois disponibles dans le
fichier cg130823a001-fra.csv. On représente cette série chronologique
par la variable aléatoire $Y_t$. L'an passé, les Canadiens ont dépensé
en moyenne $674\$$ en cadeaux au mois de décembre 2012. Notre objectif
est de prévoir quel sera le montant dépensé pour l'achat de cadeaux en
décembre 2013.

\begin{enumerate}
\item Tracez un graphique de la série chronologique $Y_t$ à l'aide
  d'un logiciel statistique. Êtes vous en mesure de déceler
  visuellement la présence d'une tendance et/ou d'une saisonnalité ?

\item Utilisez l'opérateur différentiel $\nabla_{12}$ afin d'éliminer
  la saisonnalité annuelle de la série chronologique $Y_t$ et obtenir
  la série $Z_t$. Tracez à nouveau un graphique avec les données
  obtenues. Remarquez-vous toujours la présence de saisonnalité ?
  Tracez le graphique de la composante de saisonnalité $s_t$.

\item Maintenant, nous voulons déceler s'il y a présence d'une
  tendance dans les données. En utilisant la méthode de la moyenne
  mobile avec $q=1$ et $q=5$, du lissage exponentiel double avec
  $\alpha=5\%$ et de la régression linéaire simple, estimer la
  tendance $\hat{m}_t$. Faire le graphique superposé des 5 tendances.

\item En utilisant le résultat de la régression linéaire précédente,
  prévoir la valeur non saisonnalisée en décembre 2013. En évaluant la
  moyenne des différences entre la série $Y_t$ et la valeur de la
  régression pour les mois de décembre des années précédente, on peut
  estimer la valeur $\hat{s}_{12}$. Ajouter cette valeur au résultat
  obtenu pour obtenir une estimation du taux d'inflation en décembre
  2013.

\item En applicant ce taux d'inflation à la donnée du problème,
  prédire le montant dépensé pour l'achat de cadeaux en décembre 2013.
\end{enumerate}

\clearpage
\subsection{Incendies}

On a estimé la tendance d’un ensemble de données d’incendie pour une
année. Cependant, suite à un problème informatique, certaines données
sont manquantes. Identifiez ces données.\\

\begin{tabular}{|l|l|l|l|}
  \hline
  \multicolumn{1}{|l|}{Mois} & \multicolumn{1}{l|}{Incendies} & \multicolumn{1}{l|}{Moyenne Mobile} &  \\ \hline
  1 & 4 & \multicolumn{1}{l|}{-} &  \\ \hline
  2 & 3 & \multicolumn{1}{l|}{-} &  \\ \hline
  3 & \multicolumn{1}{l|}{a} & 4,8 &  \\ \hline
  4 & \multicolumn{1}{l|}{b} & 4,8 &  \\ \hline
  5 & 2 & 5,4 &  \\ \hline
  6 & 4 & 5,2 &  \\ \hline
  7 & 6 & 3,6 &  \\ \hline
  8 & \multicolumn{1}{l|}{c} & 3,6 &  \\ \hline
  9 & \multicolumn{1}{l|}{0} & 4,4 &  \\ \hline
  10 & 2 & 3,8 &  \\ \hline
  11 & 8 & \multicolumn{1}{l|}{-} &  \\ \hline
  12 & 3 & \multicolumn{1}{l|}{-} &  \\ \hline
\end{tabular}

\clearpage
\subsection{Option de vente}

Nous sommes le 20 juin 2013, à l'heure de la fermeture des marchés
financiers. Vous possédez un titre de la compagnie BlackBerry dont la
valeur est de $S_0 = 13,82$. Un analyste vous suggère d'acheter une
option de vente européenne d'échéance 3 mois avec un prix d'exercice
équivalant à la valeur d'un contrat à terme de même échéance afin de
couvrir le risque de baisse de la valeur de ce titre. On considère des
rendements mensuels et un taux sans risque composé continument de
1,75\%.  En utilisant les logarithmes des valeurs historique à la
fermeture du titre (prendre le jour précédent si le 20 du mois n'est
pas un jour ouvré) disponibles dans le fichier \emph{blackberry.csv},
ainsi que la méthode de différenciation, évaluez la composante de
tendance $m_t$ qui équivaut aux rendements du titre. Ensuite, en
évaluant la moyenne et l'écart-type de cette composante, il est
possible d'estimer la tendance linéaire $\mu$ et la volatilité
$\sigma$ mensuelles de la série des rendements. En supposant que le
prix peut prendre deux valeurs après 3 mois, soit $S_0u = S_0
e^{3(\mu+\sigma^2/2)}$ et $S_0d = S_0 e^{3(\mu-\sigma^2/2)}$, évaluez
le prix de l'option de vente en utilisant la probabilité neutre au
risque $p^{*} = \frac{e^{rt}-d}{u-d}$ et évaluez le profit que vous
effectuerez en exercant l'option de vente le 20 septembre 2013,
considérant que vous empruntez au taux $r+2\%$ pour acheter l'option.

\clearpage

\includegraphics[height=7mm,keepaspectratio=true]{by-sa}\\%
Cette création est mise à disposition selon le contrat
\href{http://creativecommons.org/licenses/by-sa/2.5/ca/deed.fr}{%
  Paternité-Partage à l'identique 2.5 Canada} de Creative Commons
disponible à l'adresse \\
http://creativecommons.org/licenses/by-sa/2.5/ca/deed.fr \\
En vertu de ce contrat, vous êtes libre de:
\begin{itemize}
\item \textbf{partager} --- reproduire, distribuer et communiquer
  l'{\oe}uvre;
\item \textbf{remixer} --- adapter l'{\oe}uvre;
\item utiliser cette {\oe}uvre à des fins commerciales.
\end{itemize}
Selon les conditions suivantes:

  \begin{tabularx}{\linewidth}{@{}lX@{}}
    \raisebox{-9mm}[0mm][13mm]{%
      \includegraphics[height=11mm,keepaspectratio=true]{by}} &
    \textbf{Attribution} --- Vous devez attribuer l'{\oe}uvre de la
    manière indiquée par l'auteur de l'{\oe}uvre ou le titulaire des
    droits (mais pas d'une manière qui suggérerait qu'ils vous
    soutiennent ou
    approuvent votre utilisation de l'{\oe}uvre). \\
    \raisebox{-9mm}{\includegraphics[height=11mm,keepaspectratio=true]{sa}}
    & \textbf{Partage à l'identique} --- Si vous modifiez, transformez
    ou adaptez cette {\oe}uvre, vous n'avez le droit de distribuer
    votre création que sous une licence identique ou similaire à
    celle-ci.
  \end{tabularx}

\end{document}
%%% Local Variables: 
%%% mode: latex
%%% TeX-master: t
%%% End: 
